\section*{Bab \ref{ch:b3:facies-models} --- Model Fasies}

\begin{solution}{exo:b3:facies-models:1}
Dari laut ke darat: lepas pantai, transisi lepas pantai dengan pasir badai,
\omterm{def:b3:facies-models:shoreface}{muka pantai}, gisik, dataran pesisir dengan rawa. Setiap lingkungan kemudian
menumpang di atas tetangganya di arah laut: garis pantai bergeser ke arah laut, yakni regresi. Sebuah
transgresi akan menumpuk mereka dengan urutan sebaliknya, dataran pesisir di
dasar dan lumpur lepas pantai di atas, sering kali dengan permukaan erosi tempat
\omterm{def:b3:facies-models:shoreface}{muka pantai} tadi bergeser melintasinya.
\end{solution}

\begin{solution}{exo:b3:facies-models:2}
(a) Isian alur atau \omterm{def:b3:facies-models:rivers}{gosong tepi}: gerusan berkerikil di dasar, perlapisan silang,
lalu riak dan lumpur berakar. (b) \omterm{def:b3:facies-models:shoreface}{Muka pantai} progradasi atau front delta:
lumpur dengan lapisan pasir badai dan \omterm{def:b3:facies-models:hcs}{stratifikasi silang hummocky} di bawahnya,
pasir \omterm{def:b3:facies-models:shoreface}{muka pantai} berperlapisan silang dan gisik di atasnya. (c) Pasir \omterm{def:b3:facies-models:rivers}{sungai teranyam} yang bertumpuk,
gumuk pasir eolian atau gosong karbonat: set perlapisan silang dengan ukuran butir
seragam, dengan foreset gumuk besar untuk kasus eolian, atau butir skeletal dan oolitik
untuk gosongnya.
\end{solution}

\begin{solution}{exo:b3:facies-models:3}
$\phi = (2.65 - 2.40)/(2.65 - 1.0) = 0.152$; dengan air asin,
$0.25/1.55 = 0.161$.
\end{solution}

\begin{solution}{exo:b3:facies-models:4}
$e^{-2\pi z/\lambda} = e^{-\pi} = 0.043$ pada $z = \lambda/2$: sekitar
\qty{38}{m} dan \qty{55}{m}. Pada \qty{20}{m}: $e^{-2\pi\times20/75} = 0.19$
dan $e^{-2\pi\times20/110} = 0.32$ dari gerak permukaan.
\end{solution}

\begin{solution}{exo:b3:facies-models:5}
$(60 - 47.6)/(189 - 47.6) = 0.088$; $(60 - 52.6)/(189 - 52.6) = 0.054$. Sebuah nilai matriks
yang tidak pasti sebesar \qty{5}{\micro\second/ft} mengubah porositas
hingga sepertiga: persamaan rata-rata waktu memerlukan nilai matriks yang diperiksa pada inti.
\end{solution}

\begin{solution}{exo:b3:facies-models:6}
$(6.0 + 3.5 + 8.0 + 2.5)/40 = 20/40 = 0.50$. Volume pori: $20\times0.22
\times10^6 = \qty{4.4e6}{m^3}$ per kilometer persegi.
\end{solution}

\begin{solution}{exo:b3:facies-models:7}
Dekat ngarai: hanya Ta, bagian atasnya tersapu aliran berikutnya; ke bawah
kipas: Tc (riak) dan Td (lanau) di bawah Te (lumpur). Arus kehilangan kecepatan
dan muatan kasarnya seiring jarak; hanya fraksi halus yang berpindah
sejauh \qty{30}{km}, dan diendapkan oleh aliran yang lebih lambat dan membangun riak, bukan
oleh jatuhan dari suspensi pekat.
\end{solution}

\begin{solution}{exo:b3:facies-models:8}
\omterm{def:b3:facies-models:rivers}{Gosong tepi} \omterm{def:b3:facies-models:rivers}{sungai bermeander}: lag di dasar alur yang tergerus,
pasir berperlapisan silang dari alur dalam, pasir berriak dari puncak gosong,
lumpur dataran banjir dengan akar. Suksesi ini merekam alur dari dasar
hingga tepian, sehingga alur itu berkedalaman sekitar \qty{9}{m} (sedikit lebih kecil setelah
pemadatan lumpur). Ini tidak mencakup dasar: itu permukaan erosi, dan
yang berada di bawahnya dapat berupa endapan lebih tua mana pun.
\end{solution}

\begin{solution}{exo:b3:facies-models:9}
$60/40\,000 = 0.0015$, \ang{0.09}: sebuah landaian. $300/1500 = 0.20$, \ang{11.3},
setelah puncak datar berkedalaman \qty{5}{m}: \omterm{def:b3:facies-models:ramp}{paparan bertepi}, dengan butir terumbu atau gosong
di tepinya, dan blok serta aliran butir di lereng di bawahnya.
\end{solution}

\begin{solution}{exo:b3:facies-models:10}
Sabuk membutuhkan waktu $w/v$ untuk melewati suatu tempat tetap, selama itu dasar
naik relatif terhadap permukaan laut dengan laju $A$ di bawahnya: pasir mengisi $h = A\,w/v$.
$2000\times0.2/20 = \qty{20}{m}$. Dengan $A = 0$ \omterm{def:b3:facies-models:shoreface}{muka pantai} hanya meninggalkan
ketebalan profilnya sendiri, berupa lembaran tipis, karena sabuk-sabuk itu bergeser saling melewati
tanpa agradasi. Rasio $A/v$ adalah tangen dari sudut lintasan
pada proposisi di Bab~\ref{ch:b3:sequence-stratigraphy}: ketebalan setiap fasies
adalah lebarnya dikali tangen tersebut.
\end{solution}

\begin{solution}{exo:b3:facies-models:11}
Massa jenis: $(2.65 - 2.30)/1.65 = 0.212$; sonik: $(86 - 52.6)/136.4 = 0.245$.
Dengan $\phi = 0.245$, $\rho_{\mathrm{ma}} = (2.30 - 0.245)/(1 - 0.245) =
\qty{2.72}{g/cm^3}$: lebih padat daripada kuarsa, seperti batu pasir yang tersemen oleh
kalsit (2.71) atau yang mengandung butir kalsit.
\end{solution}

\begin{solution}{exo:b3:facies-models:12}
Lumpur lepas pantai di atas batu bara membawa lingkungan yang tidak berdampingan ke dalam kontak: sebuah
\omterm{def:b3:sequence-stratigraphy:parasequence}{permukaan genang}, di sini \omterm{def:b3:sequence-stratigraphy:surfaces}{permukaan transgresif} akibat erosi gelombang, dengan cangkang
puing sebagai lag. Ketika garis pantai bergeser ke arah darat, gelombang memotong dasar
hingga kira-kira \omterm{def:b3:facies-models:shoreface}{dasar gelombang cuaca cerah}: gisik, \omterm{def:b3:facies-models:shoreface}{muka pantai} atas, dan
bagian atas dataran pesisir, hingga sekitar \qty{10}{m} penampang, mungkin
hilang.
\end{solution}

\begin{solution}{pb:b3:facies-models:1}
\textbf{1.} Kapur (dingin, massa jenis rendah), serpih dari \qty{766}{m} (panas), batu gamping dari
\qty{866}{m} (dingin, padat, cepat), batu pasir dari \qty{1046}{m} (bentuk
lonceng), sekis dari \qty{1166}{m}.

\textbf{2.} Kapur \qty{766}{m}, serpih 100, batu gamping 180, batu pasir
\qty{120}{m}; dasar lapisan penutup pada $441.0 - 1166 = \qty{-725}{m}$.

\textbf{3.} $(2.71 - 2.03)/1.71 = 0.40$ dekat permukaan, $(2.71 -
2.36)/1.71 = 0.20$ pada \qty{750}{m}: kapur kehilangan separuh porositasnya
di bawah beban lapisan di atasnya, akibat pemadatan.

\textbf{4.} Butirnya berupa lempung dengan massa jenis yang tidak pasti, dan airnya
sebagian terikat padanya: asumsi persamaan tidak berlaku.

\textbf{5.} Batuan dasar kristalin: lapisan penutup ditembus, dan tidak ada reservoar
yang diharapkan di bawahnya.

\textbf{6.} $(2.71 - 2.573)/1.71 = 0.080$.

\textbf{7.} $(58.9 - 47.6)/141.4 = 0.080$: sama.

\textbf{8.} $(2.65 - 2.32)/1.65 = 0.200$.

\textbf{9.} $(79.9 - 52.6)/136.4 = 0.200$: sama.

\textbf{10.} $0.33/1.55 = 0.213$, satu poin dan sepertiga lebih.

\textbf{11.} Batu pasir, dengan porositas \qty{20}{\percent} pada pasirnya;
serpih di atas batu gamping menyegelnya, dan lumpur dataran banjir membagi
batu pasir menjadi kompartemen-kompartemen.

\textbf{12.} Berbentuk lonceng: sinar gamma rendah pada pasir silang-lapis,
naik melalui interval menghalus hingga nilai tinggi pada lumpur berakar, pada dasar yang tajam.

\textbf{13.} \omterm{def:b3:facies-models:association}{Suksesi menghalus ke atas}: lag dasar alur pada dasar erosional,
gumuk alur (silang-lapis), lalu puncak \omterm{def:b3:facies-models:rivers}{gosong tepi} (riak), dan dataran banjir
(lumpur berakar).

\textbf{14.} Alur, \omterm{def:b3:facies-models:rivers}{gosong tepinya} pada tepi dalam, tanggul alam, dan dataran
banjir dengan tanah, berdampingan.

\textbf{15.} Kira-kira tebal pasir dan interval menghalus, $10.4 + 3.2 =
\qty{13.6}{m}$, sedikit lebih besar sebelum pemadatan.

\textbf{16.} \omterm{def:b3:facies-models:rivers}{Sungai bermeander} yang sabuk alurnya berpindah berulang kali
melalui \omterm{def:b3:facies-models:rivers}{avulsi} selama cekungan mengalami penurunan, setiap alur baru memotong
dataran banjir yang ditinggalkan alur sebelumnya.

\textbf{17.} Tidak: setiap dasar bersifat erosional; lumpur dataran banjir di
bawahnya dan pasir alur di atasnya tidak harus bertetangga.

\textbf{18.} $(60 - 30)/(105 - 30) = 0.40$.

\textbf{19.} Terendah: $14\times(0.60 + 0.40\times0.25) = \qty{9.80}{m}$;
tertebal: $16\times(0.65 + 0.40\times0.20) = \qty{11.68}{m}$.

\textbf{20.} Pasir bersih \qty{58.8}{m}; interval menghalus \qty{33.4}{m}; lumpur
\qty{27.8}{m}.

\textbf{21.} Neto $= 58.8 + 0.40\times33.4 = \qty{72.2}{m}$; \omterm{def:b3:facies-models:ntg}{rasio neto terhadap
bruto} $72.2/120 = 0.60$.

\textbf{22.} Fraksi 0.20 dan 0.60: \qty{65.5}{m} dan \qty{78.8}{m}, selisih
\qty{13}{m} hanya dari pemilihan batas cutoff.

\textbf{23.} $72.2\times0.20\times10^6 = \qty{1.44e7}{m^3}$ per kilometer
persegi.

\textbf{24.} Tebal pasir neto batu pasir Trias pada Varden-2 adalah
\textbf{\qty[text-series-to-math=true]{72.2}{m}}.
\end{solution}
